\fbox{例題 1.1}

データ

\begin{tabular}{c|c|c} \hline
  No. & $x$ & $x^2$ \\ \hline
  1 & 16.8 & 282.24 \\
  2 & 14.7 & 216.09 \\
  3 & 18.3 & 334.89 \\
  4 & 15.2 & 231.04 \\
  5 & 16.4 & 268.96 \\
  6 & 17.1 & 292.41 \\
  7 & 15.8 & 249.64 \\
  8 & 16.0 & 256.00 \\ \hline
  計 & 130.3 & 2131.27 \\ \hline
\end{tabular}

\vskip\baselineskip

\noindent 1) 平均
\begin{equation*}
  \bar{x}
  = \frac{1}{8} \sum_{i=1}^{8} x_{i}
  = \frac{16.8 + 14.7 + \cdots + 16.0}{8}
  = \frac{130.3}{8}
  = 16.29
\end{equation*}

\vskip\baselineskip
\noindent 2) メディアン

データを大きさの順に並べると
\begin{equation*}
  14.3, 15.2, 15.8, 16.0, 16.4, 16.8, 17.1, 18.3
\end{equation*}
となり、$n$は偶数だから中央の2つの値、すなわち16.0と16.4の平均を求めて
\begin{equation*}
  Me = \frac{16.0 + 16.4}{2} = 16.2
\end{equation*}

\vskip\baselineskip
\noindent 3) 平方和
\begin{align*}
  S = &\sum_{i=1}^{8} {x_{i}}^2 - \frac{\left( \sum_{i=1}^{8} x_{i} \right)^{2}}{8} \\
  = &\frac{16.8^{2} + 14.7^{2} + \cdots + 16.0^{2}}{8} - \frac{130.8^{2}}{8} \\
  = &2131.27 - 2122.26 = 9.01
\end{align*}

\vskip\baselineskip
\noindent 4) 分散
\begin{equation*}
  V
  = \frac{S}{n - 1}
  = \frac{9.01}{7}
  = 1.287
\end{equation*}

\vskip\baselineskip
\noindent 5) 標準偏差
\begin{equation*}
  s
  = \sqrt{V}
  = \sqrt{1.287}
  = 1.13
\end{equation*}

\vskip\baselineskip
\noindent 6) 範囲
\begin{equation*}
  R
  = x_{text{max}} - x_{\text{min}}
  = 18.3 - 14.7
  = 3.6
\end{equation*}
